\section{Constructing the branching block}\label{sec:block}

We now construct a building block with three boundary components.  Its
essential property is that separately constant voltages produce constant
flux densities in prescribed torus coordinates.  All constants in this
section may depend on this one fixed block.

\begin{proposition}[Branching block]\label{prop:branching-block}
Let $B\subset\R^3$ be a bounded connected Lipschitz domain whose boundary
has three components $\Gamma_0,\Gamma_1,\Gamma_2$.  Suppose these components
have disjoint product collars parametrized by
$\T^2\times[0,\eta_i]$, where $\T^2=(\R/2\pi\mathbb Z)^2$.
Assume the parametrizations are bi-Lipschitz, smooth with nonsingular
derivative on open pieces covering up to a null set, and that removing
sufficiently thin collars leaves a connected Lipschitz central region.
Let $\mu_i$ be the probability measure obtained by pushing forward
$(2\pi)^{-2}\dd\theta$ to $\Gamma_i$ by its collar parametrization.

There is a scalar conductivity $\gamma_B$, bounded above and bounded away
from zero, with the following property.  For every $c=(c_0,c_1,c_2)\in\R^3$,
the weak $\gamma_B$-harmonic function $U_c$ with boundary value $c_i$ on
$\Gamma_i$ satisfies
\begin{equation}\label{bl:constant-flux}
 \int_B\gamma_B\nabla U_c\cdot\nabla\psi\,\dd x
   =\sum_{i=0}^2 I_i(c)\int_{\Gamma_i}\psi\,\dd\mu_i,
 \qquad \psi\in H^1(B),
\end{equation}
where $\sum_i I_i(c)=0$.  The linear map $c\mapsto I(c)$ has kernel
$\R(1,1,1)$ and image $\{I\in\R^3:\sum_i I_i=0\}$.
\end{proposition}

We first construct a uniformly elliptic symmetric tensor and two fields
with the required Cauchy data.  Only at the end will we apply
Theorem~\ref{thm:scalarization}.

\subsection{Attaching flat ends}

Write $B_0$ for a connected Lipschitz central region after removing
collars.  Attach an auxiliary cylinder $\T^2\times[0,\infty)$ to each
component of $\partial B_0$, using the corresponding angular coordinates.
The coordinate $t$ increases away from $B_0$.  These infinite cylinders
are auxiliary: we will retain only finite portions of them.

On $B_0$ use the identity tensor.  On each cylinder use the energy
\begin{equation}\label{bl:flat-energy}
 \int_0^\infty\int_{\T^2}
 \bigl(\abs{\partial_t p}^2+
       \nabla_\theta p\cdot S\nabla_\theta p\bigr)
 \,\dd\theta\,\dd t,
\end{equation}
where $S$ is a fixed positive definite symmetric $2\times2$ matrix such
that the numbers $h^tSh$, $h\in\mathbb Z^2\setminus\{0\}$, have no
equalities except those between $h$ and $-h$.  Such an $S$ exists:
each unwanted equality cuts out a proper hyperplane in the space of
symmetric matrices, and countably many proper hyperplanes do not cover
the open positive cone.  The hyperplane is proper because
$hh^t=ll^t$ implies $l=\pm h$.

Set $\lambda_h=(h^tSh)^{1/2}$.  For a trace $f$ on a torus, let
$\bar f=(2\pi)^{-2}\int f$ and write
$f=\bar f+\sum_{h\ne0}\widehat f(h)e^{ih\cdot\theta}$.
Ellipticity gives $\lambda_h\ge c_S\abs h$ for some $c_S>0$, so only
finitely many frequencies have rate below any fixed bound.
Given a slope $s$, its continuation into the end is
\begin{equation}\label{bl:flat-extension}
 p(t,\theta)=\bar f+st+
       \sum_{h\ne0}\widehat f(h)e^{-\lambda_h t}e^{ih\cdot\theta}.
\end{equation}
Only the decaying part has finite energy on the infinite end; the full
field is locally $H^1$ there.  Its derivative at $t=0$, directed into the
end, is $s-Pf$, where $P$ is the Fourier multiplier $\lambda_h$ on
nonconstant modes and zero on constants.

For prescribed slopes $s_0+s_1+s_2=0$, solve on $B_0$, modulo a common
constant,
\begin{equation}\label{bl:central-variational}
 \int_{B_0}\nabla p\cdot\nabla\psi\,\dd x
 +\sum_{i=0}^2\langle Pp_i,\psi_i\rangle
   =\sum_{i=0}^2s_i\int_{\T^2}\psi_i\,\dd\theta.
\end{equation}
Here $p_i,\psi_i\in H^{1/2}(\T^2)$ are traces in the attachment
coordinates.  The brackets denote the
$H^{-1/2}(\T^2),H^{1/2}(\T^2)$ duality, equivalently
\[
 \langle Pf,g\rangle
   =(2\pi)^2\sum_{h\ne0}\lambda_h\widehat f(h)\widehat g(-h).
\]
The trace forms on the left are continuous on $H^1(B_0)$.  For
completeness, in a short collar the estimate for a Fourier mode with
rate $m>0$ is
\[
 m\abs{q(0)}^2\le C\int_0^\eta
       \bigl(\abs{q'}^2+(1+m^2)\abs q^2\bigr)\,\dd t.
\]
It follows by applying the fundamental theorem of calculus and
Cauchy--Schwarz on an interval of length comparable to
$\min\{\eta,m^{-1}\}$; summation over modes and the bi-Lipschitz collar
bounds give the assertion.  The central Dirichlet energy is coercive
modulo constants by the Poincar\'e inequality on connected $B_0$; the added
trace forms are nonnegative.  The right side kills constants.  Hilbert
space minimization therefore gives Equation~\eqref{bl:central-variational}.
Together with Equation~\eqref{bl:flat-extension}, it gives a joint weak
solution, since the central outward flux functional equals $s_i-Pp_i$.

Choose two fields corresponding to a basis of
$\{s\in\R^3:\sum_i s_i=0\}$.  On each end the slope functional on this
two-dimensional space is nonzero.  Make an invertible affine change of
these two fields, on this end only, so that the new fields $u,v$ have
respective slopes $1,0$ and both have zero constant offset.  At the end
of our modification we will undo this same affine change.  Consequently
the original global pair and its slope vectors will be preserved.

If $v$ is not zero on the end, its first nonzero Fourier eigenspace
consists of a single pair of opposite frequencies.  A unimodular change
of torus coordinates, followed by a translation, puts that leading term
in the form
\begin{equation}\label{bl:leading-mode}
 b e^{-\lambda t}\cos(kx),\qquad b\ne0,\quad k\in\N,
 \quad\lambda>0.
\end{equation}
Indeed, divide the frequency vector by the greatest common divisor of
its coordinates and complete the resulting primitive vector to an
integral unimodular basis.  Translation absorbs the phase.  The angular
tensor remains constant and positive definite under this change.
Both changes preserve uniform probability measure on the torus, so they
do not alter the measures prescribed in Proposition~\ref{prop:branching-block}.

\subsection{Matching one flat end exactly}

The next lemma replaces the asymptotic form of the two fields by exact
data on a finite section of the end.  This leaves a single mode for the
frequency construction in the following subsection.

\begin{lemma}[Exact matching on one flat end]\label{bl:end-matching}
Fix one end after the preceding affine and angular normalizations, and
write $A_0=\diag(1,S)$ in its coordinates $(t,x,y)$.  Let
$u_{\rm in},v_{\rm in}$ be the two flat continuations in
Equation~\eqref{bl:flat-extension}, with zero constant offsets and
respective slopes $1,0$.  Suppose first that $v_{\rm in}$ is nonzero,
with leading term $b e^{-\lambda t}\cos(kx)$ as in
Equation~\eqref{bl:leading-mode}; thus $\lambda^2=k^2S_{xx}$.

For every sufficiently large $T$, there are smooth fields $u_T,v_T$ on
$\T^2\times[T,T+4]$ and a smooth symmetric matrix $H_T$ compactly
supported in the interior of this band such that
\begin{equation}\label{bl:end-matching-equations}
 \begin{gathered}
 \Div\bigl((A_0+H_T)\nabla u_T\bigr)=
 \Div\bigl((A_0+H_T)\nabla v_T\bigr)=0,\\
 \norm{H_T}_{L^\infty}\longrightarrow0\qquad(T\to\infty).
 \end{gathered}
\end{equation}
Near the left edge the pair is exactly $(u_{\rm in},v_{\rm in})$;
near the right edge it is exactly $(t,b e^{-\lambda t}\cos(kx))$.
The fluxes in the positive $t$ direction there are respectively
$(\partial_tu_{\rm in},\partial_tv_{\rm in})$ and
$(1,-\lambda b e^{-\lambda t}\cos(kx))$.  For large $T$, the corrected
tensor $A_0+H_T$ is uniformly elliptic.  The construction also has
$\partial_tu_T>0$, and the two gradients are independent off a union
of $2k$ smooth walls.  If $v_{\rm in}$ is identically zero, one instead
obtains $v_T=0$, right-hand field and flux pairs $(t,0)$ and $(1,0)$,
and all the stated conclusions for $H_T,u_T$.
\end{lemma}

The gradients of the target pair become dependent on the $2k$ walls
$\sin(kx)=0$.  A correction based only on the potential coordinates
$(u,v)$ therefore cannot cover the band.  The proof will remove the
source near each wall, use wall fluxes to balance the remaining source
between the regular strips, and then solve separately in those strips.
We first record the compatibility conditions and the regular solve.

All changes of coordinates in this argument are made in the energy
form.  Thus matrices, their flux columns, and sources in a new chart
are the corresponding matrix, vector, and source \emph{densities}.
The divergence equations, symmetry of a matrix density, and integrals
of source densities are preserved by these changes of coordinates.

For a compactly supported symmetric correction $H$, the two sources
$\Div(H\nabla u)$ and $\Div(H\nabla v)$ have zero integral.  Symmetry
also gives
\begin{equation}\label{bl:torque-divergence}
 v\Div(H\nabla u)-u\Div(H\nabla v)
   =\Div\bigl(vH\nabla u-uH\nabla v\bigr),
\end{equation}
so their cross-moment has zero integral as well.  On a connected region
where the gradients are independent, these conditions are sufficient
for data of fixed compact support, as the following lemma shows.

First recall an elementary compact divergence construction.  On a
rectangular box $Q$, every $r\in C_c^\infty(Q)$ with $\int_Qr=0$ is the
divergence of a vector in $C_c^\infty(Q;\R^d)$.  To see this inductively,
write $Q=(a,b)\times Q'$, choose $\eta\in C_c^\infty(a,b)$ with
$\int\eta=1$, and put $R(x')=\int_a^b r(t,x')\,\dd t$.  The first
component
\[
 F_1(x)=\int_a^{x_1}\bigl(r(t,x')-\eta(t)R(x')\bigr)\,\dd t
\]
is compactly supported, and the remaining source is $\eta(x_1)R(x')$.
Apply the construction in $Q'$ and multiply its components by $\eta$.
The one-dimensional case is direct integration.  With fixed boxes and
fixed auxiliary bumps, this is a linear construction with bounds in
each smooth norm.  The same assertion holds for data supported in a
fixed compact subset of a connected open region: use a finite connected
cover by boxes, a partition of unity, and bumps in overlaps to transfer
the mean of each piece along a spanning tree.  The final piece has zero
mean because the original data do.

\Needspace{15\baselineskip}
\begin{lemma}[Two symmetric divergence columns]\label{bl:regular-correction}
Suppose $u,v$ are smooth with independent gradients on a connected open
region.  Let $r^1,r^2$ be smooth densities supported in a fixed compact
subset of that region and satisfying
\begin{equation}\label{bl:three-moments}
 \int r^1=\int r^2=
 \int(vr^1-ur^2)=0.
\end{equation}
Then there is a compactly supported smooth symmetric matrix density $H$
such that
\begin{equation}\label{bl:correction-equations}
 \Div(H\nabla u)=r^1,\qquad \Div(H\nabla v)=r^2.
\end{equation}
For a fixed finite system of regular coordinate charts and supports,
the construction has bounds in smooth norms.  These bounds persist
under sufficiently small smooth perturbations of the charts and fields.
\end{lemma}

\begin{proof}
Use local coordinates $y_1=u,y_2=v,y_3$ and cover the source support by
finitely many such coordinate boxes joined by overlaps in the regular
region.  A partition of unity splits the data.  In any nonempty open
overlap, smooth test pairs generate arbitrary vectors of the three
moments in Equation~\eqref{bl:three-moments}.  For example, take two bumps
in the first component with different averages of $v$, and one bump in
the second component.  Their moment vectors are
$(1,0,\bar v_1)$, $(1,0,\bar v_2)$, and $(0,1,-\bar u_3)$, which are
independent when $\bar v_1\ne\bar v_2$.  Such bumps exist because $v$
is a coordinate.  Transfer the moment vector of each piece along a
spanning tree to one final box; the total vector there is zero.  We have
reduced to a pair satisfying all three conditions in a single box.

Solve $\Div F^i=r^i$ compactly by the box construction.  Put
$q=F^1_2-F^2_1$.  Integration by parts gives
\[
 \int(y_2r^1-y_1r^2)=-\int q=0.
\]
Choose a compactly supported vector $h$ with $\Div h=q$.  Add to the
first column the vector
\[
 (\partial_2h_1,\ -\partial_1h_1-\partial_3h_3,\ \partial_2h_3)^t
\]
and to the second column the vector
\[
 (\partial_2h_2,\ -\partial_1h_2,\ 0)^t.
\]
Both additions have zero divergence, and their change to
$F^1_2-F^2_1$ is $-\Div h=-q$.  The corrected columns can therefore
be completed to a symmetric matrix: set its third-row counterparts by
symmetry and its unused $33$ entry to zero.  Since
$\nabla y_1=e_1$ and $\nabla y_2=e_2$, its first two columns give
Equation~\eqref{bl:correction-equations}.  Transform back to the original
coordinates.

All operations involve fixed compact supports, fixed integrations,
smooth multipliers, and finitely many derivatives.  The overlap moment
matrices have nonzero determinants; they and the coordinate Jacobians
remain invertible under small smooth perturbations.  This proves the
stated uniform bounds, allowing a fixed finite loss of derivatives.
\end{proof}

\begin{proof}[Proof of Lemma~\ref{bl:end-matching}]
We first treat the nonzero case.  Write the translated band as
$0<t<4$, and normalize the input fields there by
\[
 \widetilde u_T(t,\theta)=u_{\rm in}(T+t,\theta)-T,
 \qquad
 \widetilde v_T(t,\theta)=\frac{v_{\rm in}(T+t,\theta)}{b e^{-\lambda T}}.
\]
They tend in every fixed smooth norm on the band to
\begin{equation}\label{bl:limiting-pair}
 U=t,\qquad V=e^{-\lambda t}\cos(kx).
\end{equation}
Indeed, the rate spectrum is discrete, so the first omitted rate in
$v_{\rm in}$, if present, is strictly larger than $\lambda$; if there
is no omitted mode, the normalized second field is already $V$.
For an $H^{1/2}$ trace, Cauchy--Schwarz bounds every differentiated Fourier series on a
translated band by the trace norm times a convergent sum of polynomial
weights against the exponential factors.  The positive rate gap then
controls the normalized remainder; the decaying part of
$\widetilde u_T$ is handled in the same way.

Choose a fixed smooth cutoff which is zero near $t=0$ and one near
$t=4$.  Interpolate from $(\widetilde u_T,\widetilde v_T)$ to $(U,V)$
with this cutoff, and call the resulting pair $u,v$.  Both endpoint
pairs are $A_0$-harmonic.  Hence the source densities
$r^1=-\Div(A_0\nabla u)$, $r^2=-\Div(A_0\nabla v)$ for the flat
tensor are supported in a fixed interior subband and tend to zero in
every fixed smooth norm.

These sources satisfy Equation~\eqref{bl:three-moments} exactly.  The
first two equalities express equality of the mean axial slopes before
and after interpolation.  For the third, integrate the counterpart of
Equation~\eqref{bl:torque-divergence} for $A_0$.  The boundary expression
is the difference of the cross-flux concomitants
\[
 \int_{\T^2}(v\partial_tu-u\partial_tv)\,\dd\theta.
\]
For the normalized input pair, continued harmonically along the
translated end, this is independent of $t$ and tends to zero at
infinity: its first field is linear plus decaying modes, and its second
has zero offset and zero slope.  For the pure pair it vanishes by angular
integration.  Hence the difference is zero.

The three global moments now vanish, but the gradients may still fail
to be independent near the walls of the limiting pair.  For
sufficiently large $T$, use coordinates $(s=u,x,y)$, since $u_t$ is
close to one.  Here $v_x$
denotes differentiation with $s,y$ held fixed.  Near each limiting
wall, the equation $v_x=0$ has a unique nearby smooth graph
$x=x_j(s,y)$ by the implicit function theorem.  Recenter by
$z=x-x_j(s,y)$.  Then
\begin{equation}\label{bl:wall-simple}
 v_z(s,0,y)=0,\qquad \abs{v_{zz}(s,0,y)}\ge c_*>0
\end{equation}
on a fixed narrow strip.  After narrowing it if necessary, $v_z/z$
extends smoothly and is bounded away from zero on the strip.  The
graphs and coordinates are periodic in $y$ and converge smoothly to
their limiting counterparts.  Choose all $2k$ strips disjoint.

We now remove the source near each wall.  In its chart, $r^1,r^2$ and
$H$ denote the transformed densities under the convention above, and
derivatives hold the other recentered coordinates fixed.  Set
\begin{align}
 \alpha(s,z,y)&=\int_0^z r^1(s,q,y)\,\dd q,\notag\\
 G(s,z,y)&=\int_0^z
       \bigl(r^2-\partial_s(\alpha v_z)\bigr)(s,q,y)\,\dd q
           -\alpha(s,z,y)v_s(s,z,y).\label{bl:wall-particular}
\end{align}
Use only the $s,z$ block of $H$, taking
\[
 H_{ss}=0,\qquad H_{sz}=H_{zs}=\alpha,\qquad H_{zz}=G/v_z.
\]
All entries involving $y$ are zero.  The quotient extends smoothly
across $z=0$: the numerator vanishes there, so $G/z$ is smooth, while
$v_z/z$ is smooth and bounded away from zero.  In particular
\[
 H_{zz}(s,0,y)=
 \frac{r^2(s,0,y)-r^1(s,0,y)v_s(s,0,y)}{v_{zz}(s,0,y)}.
\]
Differentiating Equation~\eqref{bl:wall-particular} gives
\[
 \partial_z\alpha=r^1,\qquad
 \partial_zG=r^2-\partial_s(\alpha v_z)-\partial_z(\alpha v_s).
\]
Since $u=s$, the two flux columns are
$(0,\alpha,0)^t$ and $(\alpha v_z,\alpha v_s+G,0)^t$; the displayed
identities prove both equations in Equation~\eqref{bl:correction-equations}.
The fixed denominator bound and the smooth convergence of the wall
charts show that this correction is small in every fixed smooth norm.
Integration only in $z$ preserves its support away from the axial ends.
Multiply it by a cutoff in $z$ which equals one in a narrower wall
strip and vanishes outside the chosen
strip.  After subtracting the divergences of these compact corrections,
the remaining sources are supported away from all walls.  They still
have three zero total moments by Equation~\eqref{bl:torque-divergence}.

The remaining source pair may nevertheless have nonzero moments in an
individual regular strip between successive walls.  We next transfer
these imbalances across the walls without recreating a source there.
To do so, choose any smooth $a(s,y)$ supported in a fixed middle axial
interval, put $H_{sz}=H_{zs}=a$, $H_{ss}=0$, and set
\begin{equation}\label{bl:wall-homogeneous}
 H_{zz}v_z=
 -2a\bigl(v_s(s,z,y)-v_s(s,0,y)\bigr)
 -a_s\bigl(v(s,z,y)-v(s,0,y)\bigr).
\end{equation}
The first divergence is zero because $a$ is independent of $z$.  For
the second, the $z$ derivative of the right side is
$-2a v_{sz}-a_s v_z$, which cancels
$\partial_z(a v_s)+\partial_s(a v_z)$.
Differentiating $v_z(s,0,y)=0$ with the recentered $z$ fixed gives
$v_{sz}(s,0,y)=0$.  Hence both differences on the right of
Equation~\eqref{bl:wall-homogeneous} are of order $z^2$.
Thus division by $v_z$ is again smooth and $H_{zz}v_z$ vanishes on
the wall.  Cut this matrix off away from the wall.  Its new sources
are confined to the adjacent regular strips.  Their three moments on
one side, up to a common orientation sign, are
\begin{equation}\label{bl:wall-moment-map}
 \int a(s,y)\bigl(1,v_s,v-sv_s\bigr)(s,0,y)\,\dd s\,\dd y;
\end{equation}
on the other side they are their negatives.  This follows by integrating
the two fluxes and Equation~\eqref{bl:torque-divergence}; their normal
values on the wall are $a$, $av_s$, and $a(v-sv_s)$.  They are flux
densities in this chart, so the wall integral uses $\dd s\,\dd y$.

The map in Equation~\eqref{bl:wall-moment-map} has a bounded right inverse
uniformly for large $T$.  On a limiting wall write
$V(s,0,y)=\sigma e^{-\lambda s}$, where $\sigma\in\{1,-1\}$, and set
\[
 w(s)=\bigl(1,-\sigma\lambda e^{-\lambda s},
                    \sigma(1+\lambda s)e^{-\lambda s}\bigr).
\]
Its three component functions are linearly independent on every open
interval: after multiplying a linear relation by $e^{\lambda s}$,
two derivatives force its exponential coefficient to vanish, and the
remaining affine coefficients then vanish.  Choose a nonnegative smooth
middle-supported function $\chi$, positive on an interval, and use
$a_j(s,y)=\chi(s)w_j(s)$, $j=1,2,3$.  The limiting moment matrix is
$2\pi\int\chi(s)w(s)w(s)^t\,\dd s$, which is positive definite.
Its determinant remains nonzero for the smoothly perturbed wall
coordinates and fields.  This provides a fixed finite-dimensional
right inverse with the asserted support and bounds.

The regular strips form a cyclic adjacency graph with $2k$ vertices.
Choose a spanning tree and transfer each strip's three-vector imbalance
towards a root.  The root imbalance is zero because the total is zero.
There are only finitely many transfers, so all corrections remain small
with the sources.  Their cutoff regions lie in fixed cores of the
regular strips.  Each regular strip is connected; away from the narrow
wall neighborhoods $v_x$ has a fixed nonzero sign and $(u,v,y)$ are
regular local coordinates.  The remaining source pair in each strip
now has all three moments zero and support in a fixed compact subset
of that strip.  Lemma~\ref{bl:regular-correction} removes it there.

Choose the wall strips and their transverse cutoffs first for the
limiting pair in Equation~\eqref{bl:limiting-pair}, with the axial
supports a positive distance from the band ends.  In the complementary
regular strips choose the source cores a positive distance from their
bounding walls and the band ends, and fix the associated finite chart
covers, overlap bumps, and right inverses.  These choices persist for
the smoothly close fields.  Transforming each local matrix density
back to the product coordinates and summing gives a correction $H$
with $\Div(H\nabla u)=r^1$ and $\Div(H\nabla v)=r^2$.  Thus both fields
solve the equations for $A_0+H$.  The fixed bounds, including the
finite derivative loss in Lemma~\ref{bl:regular-correction}, and the
convergence of the sources in every smooth norm give
$\norm{H}_{L^\infty}\to0$ as $T\to\infty$.  These constants are for the fixed
end; no estimate uniform in $\lambda$ or $k$ is used.  For large enough
$T$, $A_0+H$ remains uniformly positive definite.

Return to the original band by setting
\[
 u_T(T+t,\theta)=T+u(t,\theta),\qquad
 v_T(T+t,\theta)=b e^{-\lambda T}v(t,\theta),
\]
and translating $H$ to $H_T$.  Linearity preserves both equations.
Near the left and right edges, respectively, the pair is the original
input and the exact target, and $H_T=0$.  The normal row of $A_0$ is
$(1,0,0)$, so the stated positive-$t$ fluxes follow there.  Finally,
$\partial_tu_T>0$ by smooth closeness to $U=t$; away from the $2k$
wall graphs, $v_x\ne0$ in the $(s,x,y)$ coordinates, so the two
gradients are independent.  This proves the nonzero case.

If $v_{\rm in}$ is identically zero, normalize and interpolate only
the first field to $t$.  Its residual is small, compactly supported,
and has mean zero because the two mean slopes agree.
The compact divergence construction on the connected band gives a small
compactly supported flux correction $f$ with $\Div f$ equal to this
residual.  Since $e=\nabla u$ stays nonzero, the symmetric
matrix
\[
 H=\frac{fe^t+ef^t}{\abs e^2}
       -\frac{(f\cdot e)ee^t}{\abs e^4}
\]
satisfies $He=f$.  It is smooth, compactly supported, and small because
$e$ stays bounded away from zero on this fixed band.  Returning to the
original band gives the stated first field, flux, and ellipticity, with
$\partial_tu_T>0$ and the zero field unchanged.
\end{proof}

\subsection{Ending a single mode in finite distance}

The frequency-changing mechanism below is related to classical
nonunique-continuation constructions~\cite{Miller}; we prove all the
two-field properties needed here directly.  We next keep $u=t$ while
terminating the remaining mode.  Throughout
this construction the normal coefficient is one and the normal-angular
cross entries vanish.  Only the symmetric angular $2\times2$ block
varies.  Thus $u=t$ is always a solution with unit axial flux, and the
normal flux of the other field is simply its $t$ derivative.

For a pure mode $f(t)\cos(px)$, a diagonal angular block with $x$ entry
\begin{equation}\label{bl:pure-angular-entry}
 a_x(t)=\frac{f''(t)}{p^2f(t)}>0
\end{equation}
solves the equation.  The unused diagonal entry can be any fixed
positive number.  At the end of the matching band, replacing the old
constant angular block by this diagonal one preserves the pure mode's
equation, because its active entry is unchanged.  It also preserves
both fields and their normal fluxes at the interface.  If necessary,
reflect the entire normalized second field on this end, including its
retained input and matching region, so that the initial amplitude is
positive; include this reflection in the affine normalization that is
undone when returning to the original global pair.

Fix a sufficiently large $L$, to be chosen once below.  In a crossing
from frequency $p$ in one angular axis to frequency $2p$ in the other,
use an interval
$\abs{t-t_0}\le L/p$ and write
\begin{equation}\label{bl:crossing-profiles}
 f(t)=b e^{-3p(t-t_0)},\qquad
 g(t)=b e^{-p(t-t_0)}.
\end{equation}
Let $z=p(t-t_0)$.  Turn on the second mode by a smooth cutoff that is
zero near $z=-L$ and one for $z\ge-L/2$; turn off the first by a cutoff
that is one for $z\le L/2$ and zero near $z=L$.  Take these cutoffs with
derivatives of order $j$ bounded by fixed constants times $L^{-j}$.
The field is the sum of the two cut-off modes.  Before correction the
angular diagonal entries are $9$ and $1/4$, the squares of the ratios
of their decay rates to their frequencies.
At least one uncut positive profile remains throughout the crossing:
the incoming cutoff is one through $z=L/2$, and the outgoing cutoff
is one from $z=-L/2$.

Only one cutoff varies at a time.  On such a region, relabel the
dominant angular variable as $x$, its frequency as $\ell$, and its
amplitude as $D(t)>0$; label the weaker variable as $y$, its frequency
as $m$, and its cut-off amplitude as $J(t)$.  The residual is
$R(t)\cos(my)$.  Since the weak-to-dominant ratio is at most $e^{-L}$
on both cutoff regions, and $m/\ell\in\{2,1/2\}$,
\begin{equation}\label{bl:crossing-smallness}
 \frac{\abs R}{\ell^2D}+\frac{\abs J}{D}\le C e^{-L}.
\end{equation}
Indeed $R$ consists of a second cutoff derivative times the uncut weak
profile and twice a first cutoff derivative times its first derivative;
scaling by $z$ gives exactly this estimate.  The amplitude comparisons
are $g/f=e^{2z}\le e^{-L}$ on the turn-on region and
$f/g=e^{-2z}\le e^{-L}$ on the turn-off region.

The following symmetric angular correction cancels the residual
\emph{exactly}:
\begin{align}
 H_{xy}=H_{yx}&=\frac{2R}{D\ell m}\sin(\ell x)\sin(my),\notag\\
 H_{xx}&=\frac{R}{D\ell^2}\cos(\ell x)\cos(my)
            -\frac{2RJ}{D^2\ell^2}\sin^2(my),\notag\\
 H_{yy}&=0.\label{bl:crossing-correction}
\end{align}
To check this, an off-diagonal entry
$a\sin(\ell x)\sin(my)$ contributes
\[
 -aD\ell m\sin^2(\ell x)\cos(my)
 -aJ\ell m\cos(\ell x)\sin^2(my)
\]
to the angular divergence of the flux.  The divergence contributions
of the two unit $xx$ entries are
\begin{align*}
 \cos(\ell x)\cos(my)&\longmapsto
       -D\ell^2\cos(2\ell x)\cos(my),\\
 \sin^2(my)&\longmapsto
       -D\ell^2\cos(\ell x)\sin^2(my).
\end{align*}
Substitution of Equation~\eqref{bl:crossing-correction} leaves just
$-R\cos(my)$.  Equation~\eqref{bl:crossing-smallness} bounds every
correction entry by $C e^{-L}$.  Choose $L$ once so large that the perturbed
angular block has eigenvalues in a fixed positive interval, independently
of $p$ and $b$, and also that $2e^{-4L}<1$.  No $t$ derivative of this
correction occurs in the equation, because its normal row and column
are zero.

Between crossings, connect the pure outgoing frequency $p$ at decay
rate $p/2$ to the next incoming rate $3p$.  Write $r=-f'/f$ and choose
a smooth transition, constant near its endpoints, such that
\[
 p/2\le r\le3p,\qquad 0\le r'\le p^2/8.
\]
It fits in an interval of length $K/p$ for one fixed sufficiently
large $K$.  Starting from a positive amplitude, the solution
$f(t)=f(t_a)\exp(-\int_{t_a}^t r(q)\,\dd q)$ stays positive throughout the
connector.  Equation~\eqref{bl:pure-angular-entry} then gives
\begin{equation}\label{bl:connector-ellipticity}
 \frac18\le\frac{f''}{p^2f}
       =\frac{r^2-r'}{p^2}\le9.
\end{equation}
An initial pure connector joins the fixed starting rate $\lambda$ at
frequency $k$ to $3k$.  Choose its positive rate to vary slowly enough
that $r^2-r'$ remains positive.  This single connector has finite
length and positive finite ellipticity bounds depending on
$\lambda/k$; it need not share the later uniform constants.  All
interfaces preserve the profile and its first derivative, hence the
normal flux.  Changes in an unused angular entry likewise create no
interface source.

Iterate the crossings with frequencies $p_j=2^j k$, alternating the
angular axes.  Their lengths and connector lengths form a convergent
geometric series.  Let $b_j$ be the center amplitude of crossing $j$,
and let $I_j>0$ be the integral of the decay rate over its following
connector.  The amplitude at the left edge of crossing $j$ is
$b_j e^{3L}$, and its outgoing amplitude at the right edge is
$b_j e^{-L}$.  Matching to the next crossing therefore gives
\begin{equation}\label{bl:amplitude-decay}
 b_{j+1}=b_j e^{-4L-I_j}\le b_j e^{-4L}.
\end{equation}
Since $p_j=2^jk$, this gives the explicit decay
\[
 p_jb_j\le kb_0(2e^{-4L})^j\longrightarrow0.
\]
On crossing $j$ and its following connector the profile is bounded by
$C b_j e^{3L}$ and its first derivatives by $C p_j b_j e^{3L}$,
with fixed constants.  The introduced mode, extrapolated backwards
from the crossing center, is never larger than the continued dominant
mode on its turn-on region.  The displayed decay of $p_jb_j$ therefore
makes the second field and its entire gradient tend uniformly to zero
at the finite limiting section $t=t_*$.  Its energy is finite, since
the gradient is bounded and the cylinder has finite length.

Extend the second field by zero for $t\ge t_*$.  To verify the weak
equation, integrate by parts through finitely many stages.  Their
interface terms cancel by the matching of normal derivatives.  At
a terminal section tending to $t_*$ the normal flux is
$\partial_t v$, which tends uniformly to zero; its contribution against
any fixed smooth test function therefore tends to zero.  The remaining
energy integrals converge by integrability.  This proves the weak
equation across the accumulation section, with no surface source.
The tensor stays bounded and uniformly positive, although its
derivatives need not be bounded near $t_*$.  Beyond that section it
can be any fixed positive angular block with normal entry one.  Add
a short final flat band if necessary.  On this band the two normalized
fields are exactly $u=t$ and $v=0$.

\subsection{Returning to the block and scalarizing}

Apply Lemma~\ref{bl:end-matching} and the ending construction to each
of the three ends.  If its second field was identically zero, only the one-field matching
step and a final flat band are needed.  Undo each end's own invertible
affine change of fields.  The two resulting global fields agree with
their original central values and fluxes and have constant terminal
voltages and constant angular flux densities.  Their integrated terminal
currents remain the two original independent slope vectors, multiplied
by the common torus area $(2\pi)^2$.  In particular those currents are
independent and each has zero sum.

Each end now has some finite length $T_i>0$.  For each $i$ separately,
map the completed cylinder to its available physical collar by a fixed
linear change from $[0,T_i]$ to the collar's transverse interval,
followed by the given collar parametrization.  Push its tensor density,
two fields, and flux pairing forward in the energy form, and retain the
central tensor and fields on $B_0$.  Each collar map is piecewise smooth
and bi-Lipschitz with finite nonsingular bounds.  These three
pushforwards on the collars, together with the central tensor, define
a bounded uniformly positive symmetric tensor $A_B$ on the physical
block.  The bi-Lipschitz constants are fixed after completing the finite
cylinders and do not depend on the crossing index.

The fields have finite central and initial collar energy, the retained
lengths are finite, and the ending construction has finite energy.
Their traces match at each attachment, so matching-trace gluing gives
two potentials in $H^1(B)$.  For every $\psi\in H^1(B)$, restrict the
test to $B_0$ and the three collars and pull back the collar restrictions.
The energy and trace pairings are continuous on the respective $H^1$
spaces, so the weak identities on the pieces apply to these tests.
At each attachment the central and collar outward flux functionals are
opposite on the same trace.  Summing the identities therefore cancels
their terms and leaves only the terminal flux functionals.  This is the
global weak identity for the glued fields for every such test.
Uniform parameter flux density on a terminal torus becomes a constant
multiple of its specified measure $\mu_i$: the flux pushforward preserves
the flux integral against every pulled-back boundary test function,
including through the axial compression.  The glued potentials have
constant traces on the three boundary components.

We verify all rank and regularity hypotheses of
Theorem~\ref{thm:scalarization}.  The seams of the collar maps, the
finitely many central attachment interfaces, all countably many stage
interfaces, including connector interfaces, and the terminal accumulation
surfaces have volume zero.  Off these sets the tensor and fields are
smooth on open pieces.  In the central region and unchanged flat ends,
the fields are real analytic, being constant-coefficient harmonic functions.  Their
gradient wedge is either not identically zero on a connected piece,
in which case its zero set has measure zero, or is identically zero.
In the latter case, on a neighborhood where one gradient is nonzero,
write the other field as $v=F(u)$.  Its equation gives
\[
 0=\Div(A\nabla v)=F''(u)\nabla u\cdot A\nabla u,
\]
so $F$ is affine there.  The critical set of a nonconstant analytic
field has measure zero; if both fields are constant there is nothing
to check.  The elementary analytic zero-set assertion follows by
locally differentiating to the first nonzero Taylor term and applying
the one-variable zero-set fact along lines and Fubini.

In a matching band with nonzero second input, $u$ is a coordinate and
the two gradients are independent off the finitely many smooth walls, by
Equation~\eqref{bl:wall-simple} and the fixed sign of $v_x$ between them.
In a pure or crossing band, $u=t$ and the angular part of the other
field is one or two cosines on different axes.  Its angular gradient
is nonzero off a null set: with one nonzero mode its zeros lie in
finitely many walls, and with both modes present they lie in their
intersections.  Above the accumulation section the second normalized
field is zero, an allowed affine dependency.  The case with just one
matched field has a nonvanishing gradient.  These properties are
unchanged by invertible affine transformations of the pair or by the
physical collar maps.  Thus the required regular rank pieces cover
$B$ up to a null set.

Theorem~\ref{thm:scalarization} now provides a bounded positive scalar
$\gamma_B$ and two scalar-conductivity harmonic fields with exactly
the same boundary values and normal fluxes as the tensor fields.
Together with constants, their boundary values span all separately
constant data.  Indeed, if a linear combination of the two fields had
one common boundary constant, energy minimization for the positive
tensor would make it a constant field, so its integrated currents
would all be zero.  Independence of the two original slope vectors
forces that linear combination to be trivial.  Their two voltage
vectors are therefore independent modulo constants, as required.

Linearity and uniqueness for $\gamma_B$ now prove
Equation~\eqref{bl:constant-flux} for every $c\in\R^3$.  Testing with
$\psi=1$ gives $\sum_iI_i(c)=0$.  If $I(c)=0$, test instead with
$\psi=U_c$; positivity forces $\nabla U_c=0$, so connectedness of
$B$ makes all $c_i$ equal.  Conversely common constant data have zero
current.  The image therefore has dimension two and is precisely the
zero-sum current plane.  This proves Proposition~\ref{prop:branching-block}.
